
\documentclass{article}
\usepackage{colortbl}
\usepackage{makecell}
\usepackage{multirow}
\usepackage{supertabular}

\begin{document}

\newcounter{utterance}

\centering \large Interaction Transcript for game `cladder', experiment `full\_v1.5\_default', episode 771 with Qwen3.5{-}27B.
\vspace{24pt}

{ \footnotesize  \setcounter{utterance}{1}
\setlength{\tabcolsep}{0pt}
\begin{supertabular}{c@{$\;$}|p{.15\linewidth}@{}p{.15\linewidth}p{.15\linewidth}p{.15\linewidth}p{.15\linewidth}p{.15\linewidth}}
    \# & \multicolumn{2}{c}{Player} && \multicolumn{2}{c}{Game Master} \\
    \hline

    \theutterance \stepcounter{utterance}  
    & & & \multicolumn{4}{p{0.6\linewidth}}{
        \cellcolor[rgb]{0.9,0.9,0.9}{
            \makecell[{{p{\linewidth}}}]{
                \texttt{\tiny{[P1$\langle$GM]}}
                \texttt{You are an expert in causal inference. The following question is not a typical commonsense query, but rather a meticulously designed question created by a professor specializing in causal inference, intended to assess the students' mastery of the course content.} \\
\texttt{Imagine a self{-}contained, hypothetical world with only the following conditions, and without any unmentioned factors or causal relationships: Gender has a direct effect on occupation and salary. Skill has a direct effect on occupation and salary. Occupation has a direct effect on salary. Skill is unobserved. For individuals who are not male and blue{-}collar workers, the probability of high salary is 26\%. For individuals who are not male and white{-}collar workers, the probability of high salary is 61\%. For individuals who are male and blue{-}collar workers, the probability of high salary is 55\%. For individuals who are male and white{-}collar workers, the probability of high salary is 68\%. For individuals who are not male and with low skill levels, the probability of white{-}collar job is 40\%. For individuals who are not male and with high skill levels, the probability of white{-}collar job is 15\%. For individuals who are male and with low skill levels, the probability of white{-}collar job is 2\%. For individuals who are male and with high skill levels, the probability of white{-}collar job is 24\%. The overall probability of high skill level is 11\%. If we disregard the mediation effect through occupation, would gender negatively affect salary?} \\
\texttt{Start your answer with 'Yes' or 'No', followed by additional reasoning or evidence to support your explanation.} \\
            }
        }
    }
     \\ \\

    \theutterance \stepcounter{utterance}  
    & \multicolumn{4}{p{0.6\linewidth}}{
        \cellcolor[rgb]{0.561,0.737,0.561}{
            \makecell[{{p{\linewidth}}}]{
                \texttt{\tiny{[P1$\rangle$GM]}}
                \texttt{No} \\
\\ 
\texttt{To determine whether gender negatively affects salary when disregarding the mediation effect through occupation, we must calculate the **direct effect** of gender on salary. In causal inference terms, this corresponds to the effect of gender on salary while holding the mediator (occupation) constant, or more formally, the controlled direct effect (CDE) or the natural direct effect (NDE), depending on the specific counterfactual definition. However, the prompt asks to "disregard the mediation effect," which implies isolating the path \$Gender \textbackslash{}rightarrow Salary\$ and removing the path \$Gender \textbackslash{}rightarrow Occupation \textbackslash{}rightarrow Salary\$.} \\
\\ 
\texttt{Since **Skill** is an unobserved confounder that affects both Occupation and Salary, and also potentially correlates with Gender (though the prompt doesn't explicitly state a direct link between Gender and Skill, the conditional probabilities of Occupation given Gender and Skill imply a complex structure), we must be careful. However, the prompt provides specific conditional probabilities for Salary given Gender and Occupation (\$P(Salary|Gender, Occupation)\$).} \\
\\ 
\texttt{If we "disregard the mediation effect through occupation," we are essentially asking: *For a fixed occupation, does being male increase or decrease the probability of a high salary compared to being not male?* This is the definition of the direct effect in a stratified analysis.} \\
\\ 
\texttt{Let's examine the data provided for each occupation stratum:} \\
\\ 
\texttt{1.  **Blue{-}collar workers:**} \\
\texttt{    *   \$P(\textbackslash{}text\{High Salary\} | \textbackslash{}text\{Not Male\}, \textbackslash{}text\{Blue{-}collar\}) = 26\textbackslash{}\%\$} \\
\texttt{    *   \$P(\textbackslash{}text\{High Salary\} | \textbackslash{}text\{Male\}, \textbackslash{}text\{Blue{-}collar\}) = 55\textbackslash{}\%\$} \\
\texttt{    *   Difference (Male {-} Not Male) = \$55\textbackslash{}\% {-} 26\textbackslash{}\% = +29\textbackslash{}\%\$} \\
\texttt{    *   Here, being male is associated with a **higher** probability of high salary.} \\
\\ 
\texttt{2.  **White{-}collar workers:**} \\
\texttt{    *   \$P(\textbackslash{}text\{High Salary\} | \textbackslash{}text\{Not Male\}, \textbackslash{}text\{White{-}collar\}) = 61\textbackslash{}\%\$} \\
\texttt{    *   \$P(\textbackslash{}text\{High Salary\} | \textbackslash{}text\{Male\}, \textbackslash{}text\{White{-}collar\}) = 68\textbackslash{}\%\$} \\
\texttt{    *   Difference (Male {-} Not Male) = \$68\textbackslash{}\% {-} 61\textbackslash{}\% = +7\textbackslash{}\%\$} \\
\texttt{    *   Here, being male is also associated with a **higher** probability of high salary.} \\
\\ 
\texttt{In both strata of the mediator (Occupation), the probability of a high salary is strictly higher for males than for non{-}males. Therefore, the direct effect of gender on salary is positive.} \\
\\ 
\texttt{The confusion often arises from the **total effect** (which includes mediation). Let's briefly check the total effect to see if Simpson's Paradox is at play (where the total effect might} \\
            }
        }
    }
    & & \\ \\

    \theutterance \stepcounter{utterance}  
    & & & \multicolumn{2}{p{0.3\linewidth}}{
        \cellcolor[rgb]{0.9,0.9,0.9}{
            \makecell[{{p{\linewidth}}}]{
                \texttt{\tiny{[GM$|$GM]}}
                \texttt{No} \\
            }
        }
    }
    & & \\ \\

    \theutterance \stepcounter{utterance}  
    & & & \multicolumn{2}{p{0.3\linewidth}}{
        \cellcolor[rgb]{0.9,0.9,0.9}{
            \makecell[{{p{\linewidth}}}]{
                \texttt{\tiny{[GM$|$GM]}}
                \texttt{no} \\
            }
        }
    }
    & & \\ \\

    \theutterance \stepcounter{utterance}  
    & & & \multicolumn{2}{p{0.3\linewidth}}{
        \cellcolor[rgb]{0.9,0.9,0.9}{
            \makecell[{{p{\linewidth}}}]{
                \texttt{\tiny{[GM$|$GM]}}
                \texttt{game\_result = WIN} \\
            }
        }
    }
    & & \\ \\

\end{supertabular}
}

\end{document}
